\documentclass[a4paper]{article}
% Español (México) con XeLaTeX: fontspec para fuentes Unicode, polyglossia
% para la división silábica y los nombres del documento en la variante mexicana.
\usepackage{fontspec}
\usepackage{polyglossia}
\setdefaultlanguage[variant=mexican]{spanish}
% Los nombres chinos van como «pinyin (original)»: xeCJK da a los caracteres
% Han su propia fuente; sin ella XeLaTeX los omite con un aviso.
% FandolSong no cubre algunos caracteres de nombres (U+73FA); WenQuanYi Zen Hei sí.
\usepackage[AutoFallBack=true]{xeCJK}
\setCJKmainfont{FandolSong-Regular.otf}
\setCJKfallbackfamilyfont{\CJKrmdefault}{WenQuanYi Zen Hei}
% polyglossia no traduce los nombres de listings.
\AtBeginDocument{\providecommand{\lstlistingname}{}\renewcommand{\lstlistingname}{Listado}%
  \providecommand{\lstlistlistingname}{}\renewcommand{\lstlistlistingname}{Índice de listados}}
\usepackage{amsmath,amssymb}
\usepackage{graphicx}
\usepackage[margin=2.35cm]{geometry}
\usepackage[most]{tcolorbox}
\usepackage{etoolbox}
\usepackage{listings}
\usepackage{xcolor}
\usepackage{hyperref}
\usepackage{booktabs,longtable,tabularx,array}
\usepackage{subcaption}
\usepackage{float}
\usepackage{enumitem}
\usepackage{microtype}
\usepackage{fancyhdr}
\usepackage{needspace}

\definecolor{kbBack}{HTML}{F2F6FA}
\definecolor{kbFrame}{HTML}{9DB4CC}
\definecolor{kbTitle}{HTML}{52708F}
\definecolor{ibBack}{HTML}{FBF5E7}
\definecolor{ibFrame}{HTML}{C9A961}
\definecolor{ibTitle}{HTML}{7D5F1E}
\definecolor{wbBack}{HTML}{FCF2F0}
\definecolor{wbFrame}{HTML}{D6A096}
\definecolor{wbTitle}{HTML}{9E4F43}
\definecolor{dbBack}{HTML}{F4F7F3}
\definecolor{dbFrame}{HTML}{AEC2AC}
\definecolor{dbTitle}{HTML}{5F7F64}

\tcbset{softbox/.style={
  enhanced,breakable,boxrule=0.8pt,arc=2pt,left=3mm,right=3mm,
  top=2mm,bottom=2mm,coltitle=white,fonttitle=\bfseries,
  toptitle=0.6mm,bottomtitle=0.6mm
}}
\newtcolorbox{knowledgebox}[1]{softbox,colback=kbBack,colframe=kbFrame,colbacktitle=kbTitle,title={#1}}
\newtcolorbox{importantbox}[1]{softbox,colback=ibBack,colframe=ibFrame,colbacktitle=ibTitle,title={#1}}
\newtcolorbox{warningbox}[1]{softbox,colback=wbBack,colframe=wbFrame,colbacktitle=wbTitle,title={#1}}
\newtcolorbox{dialoguebox}[1]{softbox,colback=dbBack,colframe=dbFrame,colbacktitle=dbTitle,title={#1}}

\lstset{
  language=Python,basicstyle=\ttfamily\small,
  keywordstyle=\color{kbTitle}\bfseries,stringstyle=\color{wbTitle},
  commentstyle=\color{dbTitle},breaklines=true,frame=single,numbers=left,
  numberstyle=\tiny\color{gray},captionpos=b,columns=fullflexible,
  keepspaces=true,showstringspaces=false,extendedchars=true
}
\BeforeBeginEnvironment{lstlisting}{\Needspace{12\baselineskip}}

\hypersetup{colorlinks=true,linkcolor=kbTitle,urlcolor=kbTitle,citecolor=wbTitle}
\setlist{itemsep=0.25em,topsep=0.4em}
\setlength{\parindent}{2em}
\setlength{\parskip}{0.25em}
\newcolumntype{L}[1]{>{\raggedright\arraybackslash}p{#1}}

\providecommand{\notetitle}{Notas del curso: [tema del curso]}
\providecommand{\noteauthors}{Elaboradas a partir de los materiales públicos de Stanford CS25}
\providecommand{\notedate}{Versión reescrita en 2026}
\providecommand{\videochannel}{Stanford Online}
\providecommand{\videopublishdate}{2022}
\providecommand{\videoduration}{[duración del video]}
\providecommand{\videourl}{}
\providecommand{\videocoverpath}{cover.jpg}
\providecommand{\coursesite}{https://web.stanford.edu/class/cs25/}
\providecommand{\repourl}{https://github.com/hqhq1025/ai-course-notes}
\providecommand{\skillversion}{wdkns/wdkns-skills@39f1a04}

\newcommand{\figsource}[1]{\par\vspace{-0.3em}{\footnotesize\color{gray}\emph{Fuente: #1}}\par}
\newcommand{\teachervoice}[1]{\begin{knowledgebox}{Nota de clase}#1\end{knowledgebox}}
\newcommand{\term}[1]{\textbf{#1}}

\pagestyle{fancy}
\fancyhf{}
\fancyfoot[C]{\thepage}
\fancyfoot[R]{\footnotesize\color{gray}\href{\repourl}{AI Course Notes}}
\renewcommand{\headrulewidth}{0pt}
\renewcommand{\footrulewidth}{0pt}

\newcommand{\makecscover}{
\begin{titlepage}
\centering
\vspace*{0.7cm}
{\Large Stanford CS25 · Transformers United\par}
\vspace{0.8cm}
{\huge\bfseries \notetitle\par}
\vspace{0.6cm}
{\large \noteauthors\par}
\vspace{0.25cm}
{\large \notedate\par}
\vspace{0.9cm}
\ifdefempty{\videocoverpath}{}{
  \includegraphics[width=0.86\textwidth,height=0.43\textheight,keepaspectratio]{\videocoverpath}\par
}
\vfill
\begin{tcolorbox}[width=0.92\textwidth,colback=black!2!white,colframe=black!45,boxrule=0.7pt,arc=2pt]
\textbf{Autor o canal del video}: \videochannel\par
\textbf{Fecha de publicación}: \videopublishdate\par
\textbf{Duración del video}: \videoduration\par
\textbf{Sitio del curso}: \href{\coursesite}{\nolinkurl{\coursesite}}\par
\ifdefempty{\videourl}{}{\textbf{Enlace del video}: \href{\videourl}{\nolinkurl{\videourl}}\par}
\textbf{Normas de elaboración}: \skillversion\ + las normas source-first / teacher-voice / visual-QA de este repositorio
\end{tcolorbox}
\end{titlepage}
\tableofcontents
\newpage
}

\usepackage{multicol}

\renewcommand{\notetitle}{Lecture 35: diseño conjunto de RL, producto e investigación}
\renewcommand{\noteauthors}{Elaborado a partir de la conferencia oficial de Karina Nguyen (OpenAI), la grabación en 1080p y subtítulos hechos por personas}
\renewcommand{\notedate}{CS25 V5 · versión reescrita source-first de 2026}
\renewcommand{\videochannel}{Stanford Online}
\renewcommand{\videopublishdate}{Clase: 2025-04-08; publicación: 2025-04-29}
\renewcommand{\videoduration}{1:12:10}
\renewcommand{\videourl}{https://www.youtube.com/watch?v=gLwiPrwUDJ8}
\renewcommand{\videocoverpath}{cover.jpg}

\newcommand{\lecturefigure}[3]{%
\begin{figure}[H]
  \centering
  \includegraphics[width=0.97\textwidth,height=0.49\textheight,keepaspectratio]{slides-images/#1}
  \caption{#2}
  \figsource{Stanford CS25 V5 official recording, #3}
\end{figure}
}

\let\standardtableofcontents\tableofcontents
\renewcommand{\tableofcontents}{%
  \small
  \setlength{\parskip}{0pt}
  \standardtableofcontents
  \normalsize
  \setlength{\parskip}{0.25em}
}

\begin{document}
\makecscover

\section{Hilo conductor de la clase: Frontier Product Research no es «terminar la investigación y luego empaquetarla»}

El título de esta clase puede malinterpretarse fácilmente como una introducción común al aprendizaje por refuerzo (Reinforcement Learning, RL). Lo que Karina Nguyen discute en realidad es un problema de sistemas más concreto: cuando el modelo entra en la escritura, el software, la colaboración, la educación y las tareas de largo plazo, \textbf{cómo quiere el producto que actúe el modelo}, \textbf{cómo convierte la investigación esa conducta en entornos y recompensas} y \textbf{cómo permite la interfaz que el usuario saque a la luz los fallos} deben diseñarse en conjunto dentro de un mismo ciclo cerrado. Si la investigación optimiza primero un benchmark desvinculado del producto y después entrega el modelo al equipo de producto para empaquetarlo, muchos fallos reales no aparecen sino hasta el lanzamiento.

\lecturefigure{slide-01-golden-age-curiosity.jpg}{La edad de oro de «la curiosidad, la creación de conocimiento y la fabricación»: la conferencia plantea primero una orientación de valores}{00:00:49--00:01:41}

\teachervoice{00:00:49--00:01:38, Karina aclara que esto no es una lecture unidireccional, sino que espera una discusión colaborativa en el aula; su punto de partida es que cada persona puede construir algo significativo con ayuda de la IA. Este juicio de valor explica por qué lo que sigue discute a la vez creativity, agency, verification y alignment.}

\begin{importantbox}{Definición: Frontier Product Research}
En esta clase, \term{frontier product research} no significa «seguir los modelos más recientes y hacer productos». Toma las capacidades todavía inestables del modelo como material de investigación, convierte las tareas reales de los usuarios en entornos de entrenamiento y evaluaciones, y luego usa la interacción con el producto para reunir contraejemplos con rapidez. Su ciclo cerrado mínimo es
\[
\begin{aligned}
\text{product belief}&\rightarrow \text{task/env}\rightarrow \text{reward/eval}\\
&\rightarrow \text{model behavior}\rightarrow \text{user evidence}\rightarrow \text{next hypothesis}.
\end{aligned}
\]
Aquí, product belief es el principio de conducta que se espera que el producto muestre al final; debe traducirse en tareas observables y no quedarse en un eslogan.
\end{importantbox}

\begin{knowledgebox}{Límites de la evidencia en esta clase}
La primera mitad usa muchas demos de producto. Estas imágenes pueden demostrar que cierta interacción funcionaba en ese momento y mostrar cómo entiende el ponente la dirección del producto, pero no bastan por sí solas para demostrar que el modelo es confiable para todos los usuarios, distribuciones y niveles de riesgo. Lo que sigue (eval, la taxonomy de rechazos, XSTest, RL environment y reward hacking) cubre precisamente las preguntas que una demo no puede responder.
\end{knowledgebox}

\subsection{Canvas: un cambio de interfaz no es decoración, sino un cambio en la estructura de la tarea}

El primer ejemplo descompone una solicitud educativa en explicación, generación de código, ejecución y visualización. Si el sistema solo ofrece un mensaje de chat, al usuario le cuesta leer el concepto, modificar el código y observar la gráfica al mismo tiempo; Canvas eleva la «conversación» a un espacio de trabajo compartido, de modo que la salida del modelo se vuelve un objeto editable. Al leer la figura, no basta con mirar la curva Gaussian: hay que ver cómo la explicación de la izquierda, el código del centro y el resultado de la ejecución de la derecha forman un ciclo verificable.

\lecturefigure{slide-02-gaussian-canvas.jpg}{La explicación de la distribución Gaussian, el código en Python y el resultado de su ejecución se colocan en el mismo espacio de trabajo}{00:02:05--00:02:57}

Para una variable aleatoria $X\sim\mathcal{N}(\mu,\sigma^2)$, la densidad es
\[
p(x)=\frac{1}{\sqrt{2\pi}\sigma}
\exp\!\left(-\frac{(x-\mu)^2}{2\sigma^2}\right).
\]
Aquí $\mu$ controla el centro y $\sigma$ controla la escala. Que el modelo dé la fórmula no significa que la enseñanza haya terminado; los símbolos abstractos se convierten en una comprensión operable solo cuando el usuario puede modificar los parámetros, ejecutar el código y observar cómo cambia la curva. Por eso el diseño del producto cambia la pregunta de evaluación: ya no se pregunta solo si la respuesta es correcta, sino también si el código se ejecuta, si la gráfica corresponde a la explicación y si las modificaciones surten efecto de forma local.

\teachervoice{00:02:05--00:04:04, el ponente usa la enseñanza de la Gaussian y capturas de pantalla de artículos para explicar la motivación de Canvas: ChatGPT era al principio una conversational UI, pero el código, los textos largos y las preguntas de seguimiento detalladas expusieron las limitaciones de la interfaz de chat. Nota de clase: el cambio de UI proviene de la estructura de la tarea, no de un rediseño visual.}

El ejemplo de comprensión de artículos muestra además un «punto de entrada para la verificación». El usuario proporciona primero la gráfica de calibration del artículo y luego pide al modelo que la explique y que genere un informe de investigación y código. El material original, la explicación del modelo y los nuevos productos deben poder rastrearse lado a lado; de lo contrario, el sistema tiende a confundir una narración fluida con un hecho.

\lecturefigure{slide-03-paper-calibration-source.jpg}{La gráfica de calibration del artículo: la explicación del modelo debe remitir a la evidencia visual original}{00:02:57--00:03:09}

\lecturefigure{slide-04-calibration-canvas-report.jpg}{El informe de calibration generado en Canvas: se selecciona una parte del contenido y se continúa preguntando}{00:03:09--00:04:12}

\begin{knowledgebox}{Primer uso: Calibration}
\term{Calibration (calibración)} describe si la confianza del modelo coincide con su tasa real de aciertos. Si, entre todas las predicciones con confianza cercana a $q$, la proporción real de aciertos también es cercana a $q$, ese intervalo está aproximadamente calibrado. Un error frecuente es confundir calibration con accuracy: un modelo puede ser preciso pero exceder su confianza, o bien ser poco preciso pero estar calibrado en sentido probabilístico. El producto necesita mostrar por separado «cuál es la respuesta» y «cuánto debe creerse».
\end{knowledgebox}

\lecturefigure{slide-05-claude-usage-bachelor-degrees.jpg}{Gráfica de correlación de un informe de uso educativo: los datos del producto pueden generar preguntas de investigación}{00:04:12--00:04:45}

Esta gráfica de uso relaciona los usos educativos de Claude.ai con distintas carreras de licenciatura en Estados Unidos. Para leer la figura, primero se observa en el eje horizontal la proporción de cada carrera o la proporción de uso, y luego la posición de los puntos que se apartan de la línea de tendencia; la gráfica puede sugerir que las distintas disciplinas tienen patrones de adopción diferentes, pero no puede demostrar que «la formación en una carrera causa el uso» ni que «el uso causa resultados de aprendizaje». La telemetry del producto puede generar hipótesis que vale la pena investigar, pero la correlación, el sesgo de selección y el autorreporte de los usuarios deben tratarse por separado.

\begin{warningbox}{Tres filtros entre los registros del producto y una conclusión científica}
Primero, los usuarios no son una muestra aleatoria de la población; segundo, el número de clics o de generaciones no equivale al éxito en la tarea; tercero, una correlación a nivel de grupo no puede explicar directamente la causalidad individual. Frontier product research puede detectar problemas rápidamente a partir del tráfico real, pero sigue siendo necesario precisar el sampling, el outcome y el counterfactual.
\end{warningbox}

\subsection{Resumen de la sección}

La apertura de la clase ya fija la escala de toda la sesión: las capacidades del modelo, la interfaz de trabajo, la evidencia de los usuarios y los objetivos del posentrenamiento forman un solo sistema. El valor de Canvas no es «tener una ventana más», sino reunir la explicación, la edición, la ejecución y la verificación en una misma interacción; el valor de los datos del producto no es que sean verdaderos por naturaleza, sino que vuelven más concretas las preguntas de investigación.
\section{Vignettes de capacidades: el producto combina varias capacidades del modelo en un trabajo que se puede completar}

Los cinco grupos de imágenes siguientes no son una lista de funciones: muestran la \term{compositionality (composabilidad)}. Un producto útil suele invocar al mismo tiempo comprensión del lenguaje, generación de código, entorno de ejecución, renderizado visual, memoria y retroalimentación iterativa. Una capacidad que obtiene una puntuación alta en un benchmark aislado puede fallar dentro de la cadena combinada; por eso cada vignette debe leerse como un grafo de tareas de extremo a extremo.

\subsection{De la generación de código a las herramientas personales}

El ejemplo del dashboard al estilo de Stripe convierte un requisito en lenguaje natural en código de frontend y en una página interactiva. Aquí existen al menos cuatro niveles de aceptación: si se extrajeron las restricciones del requisito, si el código se ejecuta, si la jerarquía visual es legible y si, después de que el usuario hace cambios, se conserva el comportamiento existente. Si falla cualquiera de estos niveles, el producto no puede explicarse solo con la «tasa de acierto por token del código».

\lecturefigure{slide-06-stripe-dashboard.jpg}{Dashboard al estilo de Stripe generado a partir de lenguaje natural: el código y la interfaz son a la vez la salida}{00:04:41--00:05:00}

\lecturefigure{slide-07-correlation-explorer.jpg}{Explorador interactivo de correlaciones para usuarios sin formación matemática}{00:05:00--00:05:51}

El explorador de correlaciones exige además que el sistema convierta conceptos en objetos manipulables. Si el usuario arrastra la distribución de los datos pero no ve cambiar en sincronía el coeficiente de correlación, la estructura del diagrama de dispersión y los contraejemplos, la interacción es solo una «animación de demostración». Una buena eval de producto debe registrar si las variables centrales son manipulables, si la retroalimentación es inmediata, si los casos excepcionales son observables y si el usuario puede corregir, a partir de la manipulación, una idea equivocada que tenía.

\begin{importantbox}{El éxito de una tarea de extremo a extremo no es la suma de submétricas}
Sean $A_1,\dots,A_5$ los eventos de que se interprete el requisito, el código sea correcto, la ejecución tenga éxito, lo visual sea legible y el usuario complete la tarea, respectivamente. El éxito de extremo a extremo es el evento conjunto
\[
P(\text{success})=P(A_1\cap A_2\cap A_3\cap A_4\cap A_5),
\]
y no el promedio de cinco tasas de acierto. Si uno de los niveles es un umbral estricto, aunque los demás se acerquen a la puntuación máxima, el conjunto queda limitado por el eslabón más débil.
\end{importantbox}

\subsection{Imágenes, dispositivos móviles y software de costo cero}

La imagen que va de un boceto a mano a una imagen generada destaca otro tipo de interface contract: el usuario aporta una intención de baja fidelidad y el modelo completa el estilo y los detalles. Que la salida sea atractiva no significa que respete la composición, las relaciones entre objetos y la intención del autor; la evaluación debe separar aesthetic quality, instruction adherence, identity preservation y editability.

\lecturefigure{slide-08-sketch-to-image.jpg}{Del boceto a mano a la imagen generada: el modelo completa los detalles, pero el autor aún necesita controlar la dirección}{00:05:51--00:06:02}

\lecturefigure{slide-09-generative-mobile-game.jpg}{Minijuego generado al instante en el dispositivo móvil: se combinan contenido, código e interacción táctil}{00:05:52--00:06:24}

\lecturefigure{slide-10-zero-cost-software.jpg}{Demo de «la creación de software tiende a costo cero»: la interfaz móvil se genera al instante a partir de la conversación}{00:06:24--00:07:14}

Estas dos imágenes de dispositivos móviles respaldan que «el costo marginal de un prototipo disminuye», no que el software de producción ya sea gratuito. Una entrega real sigue incluyendo aclaración de requisitos, pruebas, seguridad de los datos, mantenimiento, distribución y responsabilidad. Sobre todo cuando el modelo genera software ejecutable, el entorno de ejecución debe aislarse: cuanto mayor sea la capacidad de generar código, menos se le pueden conceder al modelo por omisión permisos arbitrarios sobre la red, los archivos y las credenciales.

\teachervoice{00:04:04--00:07:36, Karina enlaza el dashboard personal, los juegos, las imágenes y el software móvil como evidencia de una «creatividad aumentada», y espera que los sistemas futuros sean más proactivos y más personalizados. Estas notas conservan su visión y añaden los límites de ingeniería: la proactividad debe diseñarse junto con los permisos, la posibilidad de revertir, el costo y la verificación.}

\begin{warningbox}{Qué es lo que más fácilmente se omite al calcular el «software de costo cero»}
El costo en tokens de generar una interfaz una vez puede ser muy bajo, pero el total cost of ownership incluye también errores en los requisitos, pruebas, infraestructura de ejecución, dependencias de terceros, auditoría de permisos, mantenimiento a largo plazo y recuperación ante incidentes. La investigación de producto debe medir el costo total de completar tareas reales, no reportar solo el precio de una generación.
\end{warningbox}

\subsection{Resumen de la sección}

La estructura común de las vignettes de capacidades es esta: la intención del usuario pasa por el modelo y la cadena de herramientas y se convierte en un artefacto observable, editable y ejecutable. Demuestran que la nueva product surface tiene valor y también exponen superficies de falla que quedan fuera de los benchmarks. La siguiente sección empieza a formular estas superficies de falla como dos scaling paradigm y dos rutas de investigación de producto.
\section{Dos Scaling Paradigm y dos rutas de investigación de producto}

El ponente distingue primero entre next-token prediction y RL on chain of thought, y después entre «una interfaz conocida que aloja una capacidad desconocida» y «primero la convicción de producto, después entrenar al modelo para lograrla». Estos dos ejes son muy importantes: el primero responde cómo se obtiene la capacidad y el segundo, cómo entra la capacidad en el producto. Si se confunden, es fácil creer por error que un preentrenamiento más potente producirá automáticamente el comportamiento de producto correcto.

\subsection{Paradigm 1: Next-token prediction es una world-building machine}

El objetivo más común del preentrenamiento autorregresivo es
\[
\mathcal{L}_{\text{NLL}}(\theta)
=-\sum_{t=1}^{T}\log p_{\theta}(x_t\mid x_{<t}).
\]
$x_t$ es el token actual, $x_{<t}$ es el contexto previo y $\theta$ son los parámetros del modelo. Un objetivo unificado permite que el modelo absorba lenguaje, conocimiento, estilo y parte de los patrones de razonamiento; por eso el ponente lo llama world-building machine. Sin embargo, lo que recompensa es la predicción local: no define directamente cuándo termina una tarea larga ni aporta señales sobre el éxito de herramientas externas, la satisfacción del usuario o la coherencia a largo plazo.

\lecturefigure{slide-11-next-token-scaling.jpg}{Paradigm 1: next-token prediction construye un modelo del mundo en múltiples dominios}{00:07:58--00:08:39}

\begin{knowledgebox}{Por qué se amplifican los errores tempranos de token}
La distribución de generación puede escribirse como $p(x_{1:T})=\prod_t p(x_t\mid x_{<t})$. En cuanto el modelo muestrea un token que se desvía del objetivo, las distribuciones condicionales siguientes se construyen sobre una historia nueva; por ello, los personajes, los hechos y los planes de un texto largo pueden desviarse gradualmente. Este problema no significa que lo autorregresivo fracase necesariamente, sino que muestra que «alta verosimilitud local» y «éxito en tareas de largo plazo» no son la misma métrica.
\end{knowledgebox}

\subsection{Paradigm 2: RL on CoT incorpora al objetivo el resultado de las tareas largas}

El aprendizaje por refuerzo denota una trayectoria de razonamiento o de acción como $\tau=(s_0,a_0,\dots,s_T)$ y optimiza
\[
J(\theta)=\mathbb{E}_{\tau\sim\pi_{\theta}}[R(\tau)].
\]
$\pi_\theta$ es la política y $R(\tau)$ es la recompensa de toda la trayectoria. Si el reward proviene de pruebas de código, verificadores matemáticos o resultados de herramientas, el modelo obtiene una señal a nivel de tarea que va más allá de la imitación token por token. En la práctica suele añadirse una restricción respecto de una política de referencia, para evitar que la política se desvíe en exceso por perseguir la recompensa:
\[
J_{\text{reg}}(\theta)=\mathbb{E}[R(\tau)]
-\beta\,D_{\mathrm{KL}}(\pi_{\theta}\Vert\pi_{\mathrm{ref}}).
\]
$\beta$ controla el equilibrio entre mejorar el comportamiento y conservar la distribución.

\lecturefigure{slide-12-rl-cot-scaling.jpg}{Paradigm 2: optimizar tareas complejas con RL sobre trayectorias de Chain-of-Thought}{00:08:39--00:09:03}

\begin{warningbox}{CoT no es automáticamente un «proceso interno» verídico}
Chain-of-Thought puede ofrecer trayectorias de cómputo más largas, admitir llamadas a herramientas y verificaciones intermedias, pero que el texto sea legible no garantiza que sea faithful. El modelo puede dar una respuesta correcta acompañada de una explicación errónea, o aprender como fórmula el formato que prefiere la recompensa. La validación debe apoyarse primero en el outcome externo, en verificaciones ejecutables y en pruebas contrafactuales, en lugar de recompensar solo lo que «parece razonamiento».
\end{warningbox}

\lecturefigure{slide-13-create-build-make.jpg}{Create, build, and make: la ampliación de capacidades debe concretarse en obras que puedan terminarse}{00:09:03--00:09:35}

\teachervoice{00:07:58--00:09:34, en clase se llama a next-token prediction y a RL on CoT dos scaling paradigm. Lo que el docente enfatiza es el cambio en la señal de supervisión: el primero aprende una distribución y el segundo puede incorporar al entrenamiento la retroalimentación sobre el resultado de tareas complejas; esto no significa que todas las tareas largas ya estén resueltas.}

\subsection{Dos rutas: Capability-first y Product-belief-first}

La primera ruta parte de una capacidad nueva y busca un form factor que el usuario ya conozca; la segunda plantea primero una convicción de producto y después construye datos, eval y entornos para que el modelo muestre ese comportamiento. La primera sirve para explorar «qué aprendió a hacer de pronto el modelo»; la segunda, para buscar una experiencia de producto coherente a largo plazo. Los proyectos reales suelen ir y venir entre ambas.

\lecturefigure{slide-14-two-product-research-paths.jpg}{Las dos rutas del research-driven product}{00:09:34--00:10:00}

\lecturefigure{slide-15-familiar-form-unfamiliar-capability.jpg}{Ruta uno: alojar una capacidad desconocida en un form factor conocido}{00:09:43--00:10:06}

\begin{importantbox}{El triángulo Capability, Affordance y Eval}
\term{Capability} es lo que el modelo puede hacer potencialmente; \term{affordance} es cómo la interfaz permite al usuario invocar, editar y deshacer; \term{eval} es la evidencia con la que el sistema juzga que algo está terminado. Los fracasos de producto suelen deberse a un desajuste entre los tres: el modelo sabe hacerlo pero la interfaz no puede expresarlo, la interfaz puede activarlo pero el eval no mide el resultado real, o el eval optimiza una métrica sustituta distinta del objetivo del usuario.
\end{importantbox}

\subsection{Caso uno: CLIP fashion search}

CLIP proyecta imágenes y textos en un espacio de embedding compartido. Si la codificación de la imagen es $f_I(i)$ y la del texto es $f_T(q)$, la puntuación de recuperación suele usar la similitud coseno
\[
s(i,q)=\frac{f_I(i)^{\top}f_T(q)}{\lVert f_I(i)\rVert\lVert f_T(q)\rVert}.
\]
Esta capacidad desconocida se colocó dentro de una interfaz de búsqueda conocida: el usuario da un texto o una prenda objetivo y el sistema devuelve resultados visualmente similares. Al leer los tres estados consecutivos, conviene comparar cómo la intención de la query cambia el conjunto de resultados, en lugar de tomar cualquier grupo de imágenes atractivas como una búsqueda exitosa.

\lecturefigure{slide-16-clip-fashion-target.jpg}{Fashion search: la recuperación parte de la prenda objetivo y sus atributos visuales}{00:10:06--00:10:11}

\lecturefigure{slide-17-clip-fashion-yellow-search.jpg}{Fashion search: la semántica del color amarillo y del modelo de prenda cambia la distribución de resultados}{00:10:11--00:10:23}

\lecturefigure{slide-18-clip-fashion-style-search.jpg}{Fashion search: una consulta de estilo en lenguaje natural devuelve un nuevo conjunto de candidatos}{00:10:23--00:10:54}

\begin{knowledgebox}{Lectura de la figura: una demo de recuperación debe medir al menos cuatro cosas}
Primero, relevance: si los resultados corresponden a la query; segundo, diversity: si solo devuelve casi duplicados; tercero, controllability: si modificar el color, el material o la silueta cambia los resultados de forma local; cuarto, failure visibility: si el sistema permite al usuario ver por qué hay coincidencia. El embedding compartido aporta el mecanismo de ordenamiento, pero no garantiza automáticamente equidad, disponibilidad en inventario ni que los términos de estilo signifiquen lo mismo en distintas culturas.
\end{knowledgebox}

\subsection{Caso dos: 100K context y carga de archivos}

Cuando la capacidad de contexto largo entra en el producto, el form factor más intuitivo es «subir un archivo y hacer preguntas». Sin embargo, que la token window sea lo bastante grande solo indica que la entrada puede recibirse, no que el modelo haya encontrado las páginas clave, agregado correctamente las cifras o mantenido la correspondencia con las fuentes. Una tarea de análisis de un 10-K de 85 páginas debe validar al menos cuatro tipos de error: retrieval, calculation, citation y omission.

\lecturefigure{slide-19-100k-context-business-analyst.jpg}{La promesa de producto de 100K context: que Claude actúe como analista de negocios}{00:10:54--00:11:00}

\lecturefigure{slide-20-100k-context-input.jpg}{Lado de entrada: un Form 10-K de 85 páginas y una solicitud de análisis concreta}{00:11:00--00:11:13}

\lecturefigure{slide-21-100k-context-analysis.jpg}{Lado de salida: los indicadores financieros clave, las tablas y las explicaciones se organizan como análisis}{00:11:13--00:11:34}

\begin{warningbox}{Long context no equivale a long-context reliability}
La ventana de contexto es un límite de capacidad; la confiabilidad depende además del sesgo de posición, los cálculos entre páginas, el análisis de tablas, la evidencia contradictoria y las citas. Si el producto solo prueba «si se puede subir», el modelo puede entregar, justo en el caso más peligroso, un resumen fluido que omite pasivos clave. Deben construirse cuatro grupos de pruebas: respuestas localizables, agregación entre secciones, ausencia de respuesta y archivos contradictorios.
\end{warningbox}

\subsection{Caso tres: Summarizer e interfaces de incertidumbre}

La demo de Summarizer concentra en una interfaz de documento los párrafos repetidos, la calibración y la edición. La convicción de producto que expresa no es «cuanto más corto el resumen, mejor», sino permitir que el usuario vea qué oraciones se conservaron, qué contenido se repite y qué conclusiones requieren nueva verificación. Si el sistema puede generar $K$ candidatos, la coherencia entre candidatos puede usarse como indicio de incertidumbre, pero un error coherente no debe tomarse como confianza. En comparación con el caso anterior de 100K context, aquí importa más cómo el usuario revisa y modifica el resultado: el contexto largo resuelve «cuánto puede leer», y la summarizer surface resuelve «cómo comprimir la evidencia en un artifact que siga siendo rastreable».

\lecturefigure{slide-22-o1-summarizer-calibration.jpg}{Summarizer: detección de repeticiones, P(IK) e interfaz de documento editable}{00:12:41--00:13:39}

\begin{knowledgebox}{De la interfaz de producto a las muestras de evaluación}
Un flujo de trabajo real puede registrar como evidence estructurada las aceptaciones, modificaciones, eliminaciones, reintentos y saltos a citas del usuario. El evento mínimo puede escribirse como
\[
e=(\text{task},\text{model version},\text{action},\text{artifact diff},\text{outcome}).
\]
Solo cuando el outcome está definido pueden los registros convertirse en datos de entrenamiento o de eval; recopilar únicamente «dónde hizo clic el usuario» lleva fácilmente a confundir desconcierto con preferencia.
\end{knowledgebox}

\begin{lstlisting}[caption={Ciclo mínimo de validación de un prototipo Capability-first}]
task = observe_real_user_job()
prototype = expose_capability_in_familiar_surface(task)
trace = run_with_permissions_and_logging(prototype)
failures = classify(trace, by=[outcome, control, trust, cost])
eval_set = convert_failures_to_reproducible_cases(failures)
ship_only_if(eval_set, product_constraints).passes()
\end{lstlisting}

\teachervoice{00:09:34--00:13:39, el ponente usa Inter Alia, 100K context y summarizer para ilustrar la primera ruta: primero aparece una capacidad desconocida y después se busca un form factor conocido. La lección práctica es que la interfaz debe permitir al usuario entender, controlar y verificar la capacidad, en lugar de limitarse a conectar el modelo a un cuadro de entrada.}

\subsection{Resumen de la sección}

Next-token prediction aporta un conocimiento previo amplio, y RL on CoT puede incorporar al objetivo el outcome de las trayectorias; ninguno de los dos determina por sí solo el producto. La ruta Capability-first necesita un form factor adecuado, y la ruta product-belief-first necesita traducir la convicción en entrenamiento y evaluación. La lección común de CLIP, el contexto largo y summarizer es que la capa de producto debe definir de forma explícita la controlabilidad y la evidencia.
\section{Product-belief-first: escribir las creencias sobre el comportamiento en la interfaz, los datos y el postentrenamiento}

La segunda vía parte de una afirmación de producto más fuerte: ¿cómo debe trabajar el sistema junto con las personas? Este tipo de objetivo a menudo no puede deducirse de un único benchmark. El equipo necesita primero construir prototipos de interacción que hagan que los usuarios revelen «dónde se parece a un colaborador y dónde se parece a un autocompletado», y después convertir esas diferencias en ejemplos de entrenamiento, reward y eval. La belief de producto no es una preferencia estética que se impone sobre la evidencia, sino una hypothesis que espera ser operacionalizada y refutada.

\lecturefigure{slide-23-product-belief-vision.jpg}{Vía dos: partir de la product belief y la vision, y después lograr que el modelo lo haga}{00:13:39--00:14:28}

\begin{importantbox}{Cómo escribir una product belief que pueda verificarse}
Ni «el modelo debe ser más inteligente» ni «la experiencia debe ser más natural» pueden verificarse directamente. Una belief operativa debe incluir: el usuario y la tarea objetivo, los comportamientos que el modelo debe adoptar o evitar, el control que conserva el usuario, la ruta de recuperación ante fallos, la evidence de éxito y los límites de costo. Por ejemplo, «ser un collaborator de escritura» implica, como mínimo, que el modelo pueda leer la selección, proponer modificaciones locales, explicar los cambios, conservar la voz del autor y permitir deshacerlos paso a paso.
\end{importantbox}

\subsection{Language as a high-end product}

La ponente considera el language en sí como material de producto de alta gama. El lenguaje no es relleno de la interfaz: el tono, la densidad de información, la estructura, el grado de franqueza, las citas y la expresión de la incertidumbre modifican la manera en que el usuario entiende el mundo. El ejemplo del producto de noticias muestra que un mismo context block puede presentarse de formas distintas en la story page, en el carousel de la página de inicio y en el bottom sheet; los hechos subyacentes se comparten, pero la organización superficial debe adaptarse a cada escenario.

\lecturefigure{slide-24-language-high-end-product.jpg}{Language as a high-end product: la organización de la información es en sí misma comportamiento del producto}{00:14:28--00:14:36}

\lecturefigure{slide-25-context-block-multiple-surfaces.jpg}{Context block compartido: los mismos hechos se mantienen coherentes entre surfaces}{00:14:36--00:14:45}

\lecturefigure{slide-26-adaptive-presentation-surfaces.jpg}{La misma información se presenta de forma adaptativa en el carousel en línea, la página de inicio y el bottom sheet}{00:14:45--00:15:39}

Puede denotarse la representación semántica compartida como $z$, la situación del usuario como $u$ y la surface como $s$; la presentación final es
\[
y_s=g_s(z,u).
\]
$g_s$ se encarga de la extensión, el diseño y la forma de interacción, pero no debe alterar en silencio los hechos contenidos en $z$. En ingeniería, conviene separar en capas el factual core, la presentation policy y el user state; de lo contrario, una edición en una surface puede generar incoherencias en otra, o el modelo puede eliminar matices clave para adaptarse a una interfaz breve.

\begin{knowledgebox}{Primer uso: Single source of truth}
Single source of truth no exige que todas las interfaces muestren la misma frase, sino que varias surfaces hagan referencia al mismo objeto de hechos rastreable. El modelo puede resumir, reordenar y reformular, pero cada conclusión debe poder remitirse a la fuente compartida; cuando el usuario actualiza un hecho en un lugar, los demás lugares deben sincronizarse explícitamente o indicar la versión.
\end{knowledgebox}

\subsection{Del Writing IDE a la micro-personalization}

El ejemplo del Writing IDE inserta el modelo de lenguaje en el flujo de trabajo del autor, en lugar de pedirle que copie todo el texto en una ventana de chat. La verdadera dificultad no es generar una frase bonita, sino entender la selección, el contexto, la intención del autor y el alcance de los cambios. La micro-personalization va más allá y exige que el sistema recuerde «cómo trabaja este usuario», pero la personalización no debe convertirse en una manipulación invisible.

\lecturefigure{slide-27-writing-ide-gpt3.jpg}{Writing IDE: el modelo ofrece colaboración local dentro del contexto del documento}{00:15:39--00:15:50}

\lecturefigure{slide-28-language-compressed-knowledge.jpg}{El lenguaje es un artifact altamente comprimido del conocimiento humano y también material para la innovación de interfaces}{00:15:50--00:16:17}

\lecturefigure{slide-29-claude-micro-personalization.jpg}{Micro-personalization en Claude.ai: organizar los accesos posteriores según los temas del historial}{00:16:17--00:16:20}

Si el estado de personalización es $m_u$, una actualización puede escribirse como
\[
m_u' = U(m_u,e_t,\gamma),
\]
donde $e_t$ es la evidence de la interacción actual y $\gamma$ es la intensidad de retención. El producto debe responder: qué eventos pueden registrarse, cuánto tiempo se conservan, cómo puede el usuario consultarlos o eliminarlos y cómo se revierte un recuerdo erróneo. Si $U$ es completamente invisible, el modelo puede consolidar una preferencia ocasional como un perfil de largo plazo.

\begin{warningbox}{Tres errores de la personalización}
Primero, tomar los clics frecuentes como preferencias estables; segundo, trasladar información sensible entre tareas sin que el usuario lo sepa; tercero, reducir la diversidad de puntos de vista con tal de resultar «atento». La micro-personalization debe medir a la vez relevance, surprise, user control y privacy, y no solo la tasa de clics.
\end{warningbox}

\teachervoice{00:13:39--00:16:18, Karina explica, a partir del producto de noticias, el Writing IDE y el carácter comprimido del lenguaje, cómo la product belief antecede al comportamiento del modelo. El docente enfatiza: la verdadera interface innovation suele requerir entrenar al modelo para que entienda nuevos objetos de interacción, y no solo dibujar una UI alrededor de un modelo existente.}

\subsection{Virtual teammate: asincronía, varias personas y acceso a herramientas}

«AI teammate» se usa con frecuencia de forma abusiva como término de marketing. El prototipo de Claude in Slack que presenta la ponente tiene límites de comportamiento más concretos: conversación asíncrona, colaboración entre varias personas, tareas automáticas, resúmenes periódicos del canal y acceso a herramientas. Todavía no es un colega completo, pero al menos reformula teammate, de un adjetivo de personalidad, a un protocolo de trabajo observable.

\lecturefigure{slide-30-claude-slack-product-spec.jpg}{Claude in Slack: generar product spec, objetivos y OKR a partir del historial del canal}{00:16:20--00:17:50}

\begin{knowledgebox}{Términos clave: Assistant, Tool, Agent, Teammate}
\begin{tabularx}{\textwidth}{L{0.17\textwidth} X X}
\toprule
Rol & Mecanismo central & Pregunta principal de verificación \\
\midrule
Assistant & Genera respuestas a solicitudes & Si la ayuda en un solo turno es correcta y clara \\
Tool & Ejecuta operaciones acotadas & Si la entrada y la salida son deterministas y pueden auditarse \\
Agent & Mantiene estado e invoca herramientas & Permisos, plan, recuperación, condiciones de paro \\
Teammate & Coordina de forma continua en flujos con varias personas & Límites de responsabilidad, contexto compartido, traspaso y juicio social \\
\bottomrule
\end{tabularx}
Un mismo modelo puede desempeñar roles distintos; el rol lo definen el entorno y los permisos, no el tono del chat.
\end{knowledgebox}

Un teammate en un canal con varias personas debe mantener, como mínimo, un estado compartido $s_t$, un estado visible para cada persona $p_i$ y un conjunto de permisos $\mathcal{A}_i$. Las acciones ejecutables deben cumplir
\[
a_t\in \mathcal{A}(s_t,p_i,\text{policy}),
\]
en lugar de ejecutarse solo porque el modelo «considera que ayudan». Resumir el canal, escribir una spec y modificar un issue tienen niveles de riesgo distintos y deben tener umbrales de confirmación distintos.

\subsection{Entrenar al modelo para que sea un collaborator}

Canvas no queda terminado con solo conectar un editor de texto al modelo. El trabajo de entrenamiento que repasa la ponente incluye datos objetivo, uso del context, formato de documentos, editing suggestion, rewriting, formatting y output quality. Los resultados en barras de la figura indican que el equipo descompuso la «colaboración» en varias tareas repetibles y después iteró sobre los datos y el postentrenamiento; no demuestran que una puntuación global represente todos los escenarios de escritura.

\lecturefigure{slide-31-human-ai-flexible-affordances.jpg}{Human--AI collaborative affordance: la interfaz debe ampliarse según la tarea}{00:17:50--00:19:36}

\lecturefigure{slide-32-training-collaborator.jpg}{Entrenar al modelo para que sea un collaborator: varios tipos de comportamiento de edición y comparación de calidad}{00:18:21--00:19:41}

\begin{importantbox}{Contract mínimo de comportamiento de un collaborator}
\begin{enumerate}
  \item Leer el artifact actual, la selección, el objetivo del autor y el alcance de las partes que no deben modificarse;
  \item Proponer primero un plan local y después generar un diff comparable;
  \item Conservar las zonas no autorizadas y la voice del autor;
  \item Dar la fuente o la incertidumbre de los cambios que afectan a los hechos;
  \item Permitir aceptar, rechazar, modificar parcialmente y hacer un rollback completo;
  \item Convertir la retroalimentación del usuario en una eval reproducible para la siguiente iteración, en lugar de limitarse a la adaptación en línea.
\end{enumerate}
\end{importantbox}

\begin{lstlisting}[caption={Ciclo de colaboración en documentos con control del usuario}]
state = load(document, selection, user_goal, permissions)
proposal = model.plan_and_edit(state)
diff = isolate_changes(proposal, protected_regions=state.locked)
evidence = verify_facts_and_style(diff, state.sources)
decision = user.review(diff, evidence)
if decision.accepted:
    commit(diff)
else:
    record_failure(state, diff, decision.reason)
\end{lstlisting}

\teachervoice{00:17:50--00:20:25, en clase se preguntó «qué significa exactamente que el modelo sea un collaborator», y se señaló que Canvas requiere entrenamiento específico en context use, formatting, rewriting y editing suggestion. Experiencia práctica: detrás de una buena interfaz de colaboración hay datos de comportamiento y eval, no un cuadro de entrada más grande.}

\subsection{Modular tool composition: Tasks no es un solo prompt}

La pantalla de ChatGPT Tasks combina la activación programada, la generación de un texto, su edición posterior y su guardado. Aquí el sistema no solo genera texto: también necesita un scheduler, un artifact store, permisos y recuperación ante fallos. Conectar varios módulos produce nuevos state mismatch: una tarea puede ejecutarse sobre un contexto desactualizado, el resultado generado puede sobrescribir las modificaciones del usuario y la ejecución repetida puede provocar conflictos.

\lecturefigure{slide-33-chatgpt-tasks-composition.jpg}{ChatGPT Tasks: combinación modular de tareas programadas, el artifact de Canvas y la edición posterior}{00:19:41--00:20:25}

\begin{knowledgebox}{Primer uso: Idempotency y rollback}
\term{Idempotency (idempotencia)} significa que ejecutar repetidamente la misma acción no produce efectos secundarios adicionales; las tareas programadas, en particular, necesitan una idempotency key para evitar que los reintentos de red publiquen o modifiquen archivos más de una vez. \term{Rollback (reversión)} exige que el sistema pueda volver a una versión conocida. Ambos son fundamentos de los productos de agent/tool y no deben añadirse hasta después de que el modelo falle.
\end{knowledgebox}

\subsection{Resumen de la sección}

El núcleo de Product-belief-first es descomponer «cómo queremos que colabore el modelo» en objetos de interfaz, permisos, datos y eval. Los productos de lenguaje, el context compartido, la personalización, el teammate en Slack, el collaborator de Canvas y Tasks muestran lo mismo: cuanto más cerca se está del trabajo real, más se necesitan versiones, permisos, diff, rollback y control del usuario. La siguiente sección aborda el caso más sólido de esta clase: cómo convertir un model behavior difuso en una eval confiable.
\section{Behavior Evals: de «algo no está bien en el modelo» a un sistema de rechazo depurable}

El comportamiento de un modelo no es como un problema de matemáticas que tiene una sola respuesta. Los rechazos, el carácter directo, el equilibrio en temas políticos, el reconocimiento de alucinaciones, el estilo lingüístico y las reacciones emocionales conforman en conjunto el producto que percibe el usuario. La ponente pregunta primero «cómo hacer una eval confiable» y después usa la mejora del over-refusal de Claude 2.1 a Claude 3 para mostrar que el proceso correcto no es ajustar por intuición un reward total, sino construir una taxonomy, combinar fuentes de datos, rastrear el problema hasta conjuntos de entrenamiento concretos y hacer regresión con ejemplos before/after y benchmarks.

\lecturefigure{slide-34-evals-question.jpg}{Pregunta central: cómo hacer una eval confiable y cómo medir el comportamiento que realmente se necesita}{00:20:25--00:20:43}

\begin{importantbox}{Definición: una eval es un protocolo de decisión, no un archivo de preguntas}
Una eval completa contiene al menos la distribución de tareas $D$, la versión del modelo/sistema $M$, la regla de calificación $S$, las condiciones de ejecución $C$ y el umbral de aceptación $\tau$:
\[
\operatorname{Eval}(M;D,S,C)\ge \tau.
\]
Si cambian el prompt, las herramientas, el system message, el sampling, la versión del judge o la rubric humana, los resultados ya no pueden compararse directamente. Lo que vuelve confiable una eval es que sea reproducible y que pueda sustentar decisiones de producto, no la cantidad de preguntas.
\end{importantbox}

\subsection{Escribir el model behavior como un vector}

El primer grupo de traits incluye reducir el benign refusal, responder con una postura más definida y disminuir las milquetoast response, al tiempo que se conservan caveats donde hay incertidumbre. El segundo grupo de traits incluye el nuance en temas de alto riesgo, fomentar la media literacy, dar explicaciones moderadamente más largas y saber lo que uno no sabe. No forman una misma dimensión de «helpfulness»; si solo se optimiza una puntuación única, el modelo puede encubrir errores factuales con textos más largos y más complacientes.

\lecturefigure{slide-35-general-behaviors-refusal-opinions.jpg}{Vector de comportamiento I: benign refusal, carácter directo, caveat y respuestas que no siguen una plantilla}{00:20:43--00:23:23}

\lecturefigure{slide-36-general-behaviors-nuance-knowledge.jpg}{Vector de comportamiento II: nuance, media literacy, longitud y calibrated self-knowledge}{00:23:23--00:24:17}

El comportamiento de una respuesta puede representarse como
\[
\mathbf{b}(x,y)=
\begin{bmatrix}
h & s & d & n & c & v
\end{bmatrix}^{\top},
\]
donde $h$ es helpfulness, $s$ es safety, $d$ es directness, $n$ es nuance, $c$ es calibration y $v$ es voice consistency. La aceptación del producto no consiste en buscar un máximo global, sino en definir una región permitida $\mathcal{B}_{\text{task}}$ para cada dominio de riesgo. Por ejemplo, las preguntas médicas exigen caveats más fuertes, mientras que la escritura creativa no debería rechazarse automáticamente porque aparezcan palabras como «vigilancia» o «heist».

\begin{knowledgebox}{Términos clave: dimensiones comunes de la evaluación de comportamiento}
\begin{tabularx}{\textwidth}{L{0.18\textwidth} X X}
\toprule
Dimensión & Problema que resuelve & Qué no puede sustituir \\
\midrule
Helpfulness & Si hace avanzar el objetivo del usuario & No sustituye la corrección factual \\
Harmlessness & Si evita daños reales & No puede juzgar la intención por palabras clave superficiales \\
Directness & Si responde primero la pregunta central & No puede omitir las salvedades necesarias \\
Nuance & Si conserva la evidencia en conflicto y las condiciones & No puede convertirse en complacer a ambas partes sin llegar a una conclusión \\
Calibration & Si expresa un grado de confianza razonable & No sustituye la verificación externa \\
Self-knowledge & Si conoce los límites de su capacidad o de su información & No puede simularse con descargos de responsabilidad fijos \\
\bottomrule
\end{tabularx}
\end{knowledgebox}

\teachervoice{00:20:25--00:23:52, la clase no resumió «Claude 3 es mejor» con una sola puntuación, sino que enumeró uno por uno refusals, directness, nuance, balance y self-knowledge. Nota de clase: primero hay que descomponer el comportamiento del producto; solo así el equipo sabe qué aspecto mejoró y cuál se sacrificó.}

\subsection{Over-refusal: el clasificador de seguridad toma las palabras superficiales por intención}

El problema típico de Claude 2.1 es que el prompt contiene en apariencia palabras peligrosas, aunque la tarea real es inofensiva. En la figura, «kill a python process» es una pregunta de gestión de procesos; si el modelo solo asocia el riesgo por keyword, la rechazará. Más difíciles de detectar son los casos de escritura creativa, de documentos largos adjuntos y de function calling: el sistema puede producir un misdirected refusal porque algún pasaje del contexto parece dañino o porque olvida el system prompt.

\lecturefigure{slide-37-claude21-overrefusal.jpg}{Over-refusal de Claude 2.1: palabras peligrosas superficiales hacen que se rechacen solicitudes inofensivas}{00:23:52--00:25:55}

\lecturefigure{slide-38-nuanced-refusals.jpg}{Rechazos más matizados: explicar la responsabilidad, reconocer los límites y ofrecer alternativas viables}{00:25:55--00:26:27}

Sea $D_b$ el conjunto de benign prompts y $r(x)\in\{0,1\}$ el indicador de rechazo del modelo; entonces la false refusal rate es
\[
\operatorname{FRR}=\frac{1}{|D_b|}\sum_{x\in D_b}r(x).
\]
Que la FRR baje no significa automáticamente que la seguridad haya mejorado, porque el modelo también puede empezar a responder solicitudes realmente dañinas. Por eso debe reportarse al menos junto con la harmful compliance rate y estratificarse por dominio; el promedio total oculta fallas locales en function calling, creative writing o documentos largos.

\begin{warningbox}{Rechazar menos no siempre es mejor}
El objetivo de optimización debería ser «rechazar menos las solicitudes inofensivas, limitar de forma robusta los riesgos reales y ofrecer ayuda segura y explicable ante las solicitudes limítrofes». Reducir solo la FRR induce unsafe compliance; reducir solo la harmful compliance induce rechazos basados en plantillas. Ambos casos requieren hard negatives: pares de ejemplos parecidos en la superficie pero con intenciones distintas.
\end{warningbox}

\subsection{Taxonomy: primero nombrar la falla y después decidir qué reparar}

La taxonomy que enumera la ponente incluye benign over-refusal, creative-writing refusal, function-calling refusal, long-document attachment y misdirected refusal. La clasificación no sirve para redactar informes, sino para mapear «el modelo rechazó» a la capa del sistema que podría causarlo: datos de seguridad, datos de self-knowledge, tool schema, enrutamiento de long-context u olvido del system prompt.

\lecturefigure{slide-39-refusal-taxonomy.jpg}{Taxonomy de rechazos: benign, creative writing y function calling}{00:26:27--00:27:08}

\lecturefigure{slide-40-refusal-taxonomy-examples.jpg}{Ejemplos visuales de la taxonomy de rechazos: adjuntos y misdirected refusal}{00:27:08--00:28:12}

\begin{knowledgebox}{Tabla de root cause: causas distintas bajo una misma apariencia de rechazo}
\begin{tabularx}{\textwidth}{L{0.22\textwidth} X X}
\toprule
Síntoma & Posible causa raíz & Qué revisar primero \\
\midrule
Se rechaza ``kill python process'' & Datos de seguridad basados en palabras clave con demasiado peso & benign hard negatives \\
Se rechaza una historia creativa & Los ejemplos de seguridad contaminan el dominio creativo & Proporción entre dominios y pair data \\
``can't see tool'' & Datos de self-knowledge erróneos & Estado de disponibilidad de las herramientas \\
Un adjunto largo provoca el rechazo & Contenido local del documento induce un juicio global erróneo & Atribución por párrafo e intención de la consulta \\
Olvida el system prompt & Defecto de contexto o de enrutamiento & Ensamblado del prompt y trazas de regresión \\
\bottomrule
\end{tabularx}
Un mismo nombre, «refusal», puede requerir reparaciones completamente distintas; agregar primero datos genéricos de SFT suele provocar contaminación cruzada.
\end{knowledgebox}

\begin{lstlisting}[caption={Flujo mínimo de triaje para fallas de rechazo}]
case = replay(prompt, attachments, tools, system_prompt, model_version)
if case.intent_is_benign and case.refused:
    bucket = classify_refusal(case)
    sources = trace_training_and_policy_influences(bucket)
    candidate_fix = design_targeted_data_or_system_change(sources)
    run_regression(candidate_fix, benign_suite, harmful_suite, domain_suite)
else:
    preserve_or_strengthen_safety_boundary(case)
\end{lstlisting}

\teachervoice{00:23:52--00:28:35, Karina enfatiza repetidamente que el over-refusal no proviene de una sola fuente de datos; el equipo construyó una taxonomy que distingue benign, creative writing, function calling, long document y misdirected refusal. Experiencia práctica: ante un mal ejemplo, primero hay que preguntar a qué mecanismo pertenece, en lugar de modificar de inmediato el reward global.}

\subsection{Combinación de fuentes para una eval confiable}

Las slides dividen los eval prompts en fallas reales recopiladas en el product flywheel, ejemplos de usuarios curados manualmente, synthetic prompts que cubren la frontera harmless/helpful y benchmarks externos como XSTest. Cada tipo atiende un punto ciego distinto: el tráfico real tiene validez ecológica pero un muestreo sesgado; los synthetic data son controlables pero heredan el sesgo del generator; los benchmarks públicos son reproducibles pero propensos al sobreajuste.

\lecturefigure{slide-41-trusted-evals-sources.jpg}{Fuentes de una eval confiable: product flywheel, synthetic prompts y XSTest}{00:28:12--00:30:21}

Si la distribución de la $k$-ésima fuente es $D_k$, la eval combinada puede escribirse como una mixture
\[
D_{\text{eval}}=\sum_{k=1}^{K}w_kD_k,
\qquad \sum_k w_k=1.
\]
$w_k$ no es una constante natural, sino una decisión de riesgo del producto. Ponderar según la proporción del tráfico subestima los riesgos altos de la cola larga; una ponderación completamente uniforme, en cambio, puede alejarse del uso real. Lo más prudente es reportar a la vez overall, slice y worst-case, y congelar los conjuntos de regresión críticos.

\begin{knowledgebox}{Primera mención: RLAIF y synthetic preference data}
\term{RLAIF (Reinforcement Learning from AI Feedback)} usa un AI judge o preferences generadas por IA para sustituir o complementar la retroalimentación humana. Los synthetic preference data permiten cubrir con rapidez estilos, límites y casos de cola larga específicos, pero su valor depende de la diversity, de la judge calibration y de la auditoría por muestreo. No se trata de «la IA demostrando por sí misma que tiene razón», sino de un proceso de producción de datos que sigue requiriendo verificación externa.
\end{knowledgebox}

\subsection{Estrategia de reparación: depurar los datos como se depura software}

Las General Approaches incluyen limpiar o regularizar los datos, la recolección humana dirigida, generar anti-refusal synthetic data y el RL directo. Lo decisivo no es el nombre del método, sino localizar la causa raíz. A continuación, la ponente da la frase práctica más importante: \textbf{look at data like you'd debug software}. El function-calling refusal puede provenir de datos de self-knowledge; los rechazos ante documentos largos, de ejemplos del tipo «no tengo capacidad visual»; y la escritura creativa puede verse afectada por la proporción de datos de seguridad.

\lecturefigure{slide-42-refusal-general-approaches.jpg}{Rutas generales para reparar el over-refusal y resultados en WildChat/XSTest}{00:30:21--00:31:00}

\lecturefigure{slide-43-debug-data-like-software.jpg}{Revisar los datos como se depura software: rechazos distintos provienen de rutas distintas}{00:31:00--00:31:13}

\lecturefigure{slide-44-dataset-specific-refusals.jpg}{Posible acoplamiento entre los rechazos en escritura creativa y los datos de seguridad}{00:31:13--00:31:44}

\begin{importantbox}{Los cuatro objetos rastreables del data debugging}
\begin{enumerate}
  \item \textbf{Provenance}: si el ejemplo proviene de usuarios, anotadores, modelos o plantillas de políticas;
  \item \textbf{Slice}: a qué dominio de tarea, de riesgo y de idioma pertenece el ejemplo;
  \item \textbf{Influence}: qué slices de comportamiento cambian a la par después del entrenamiento;
  \item \textbf{Reversibility}: si es posible retirar esos datos o reducir su peso y volver a ejecutar experimentos locales.
\end{enumerate}
Una gran mezcla de datos sin provenance es difícil de depurar; «agregar algunos ejemplos buenos más» puede ocultar el conflicto en lugar de resolverlo.
\end{importantbox}

\teachervoice{00:28:35--00:31:44, la clase pasa de los product prompts, los synthetic prompts y XSTest a la root cause en los datos. La docente dice explícitamente que cada tipo de refusal puede deberse a un conjunto de datos distinto y que hay que revisarlo como cuando se depura software, en lugar de confiar en una perilla genérica de safety/helpfulness.}

\subsection{Helpfulness--Harmlessness es una optimización con restricciones, no un control deslizante único}

El data debugging anterior mostró que cada tipo de rechazo tiene una causa raíz distinta; ahora hay que resolver una decisión de producto que abarca todas las categorías: cuando el modelo está más dispuesto a responder, puede aumentar a la vez la información que infringe las políticas; cuando es más conservador, perjudica a los usuarios legítimos. Este problema no se resuelve con un control deslizante de «nivel de seguridad», porque el false refusal y el harmful compliance no son simétricos en costo, distribución ni frontera de responsabilidad. El objetivo puede escribirse como una optimización con restricciones:
\[
\max_{\pi}\;\mathbb{E}[H(\pi)]
\quad \text{s.t.}\quad
\mathbb{E}[U(\pi)]\le \epsilon,
\]
donde $H$ es la utilidad, $U$ es el riesgo inaceptable y $\epsilon$ es el presupuesto de riesgo que fija el producto. Un sistema real también debe establecer restricciones distintas por dominio y considerar el costo asimétrico de los false positive y los false negative.

\lecturefigure{slide-45-helpfulness-harmlessness.jpg}{Helpfulness y Harmlessness: buscar el equilibrio, no elegir un solo lado}{00:31:44--00:31:54}

\lecturefigure{slide-46-xstest-refusal-rates.jpg}{Incorrect refusal rate en XSTest: comparar versiones exige entender bien los ejemplos y la métrica}{00:31:54--00:32:10}

La figura de XSTest muestra diferencias notables en la incorrect refusal rate entre versiones de Claude. Al leer la figura, primero hay que confirmar que el eje vertical mide los rechazos erróneos y no la seguridad general, y después comparar las versiones del modelo; una barra baja respalda que hay menos rechazos ante benign prompts, pero no demuestra mayor seguridad ante harmful prompts ni representa todos los escenarios de idioma, herramientas o documentos largos. La cifra total de la figura debe interpretarse junto con las taxonomy slices.

\begin{warningbox}{Una mejora en el benchmark no significa que el problema del producto esté cerrado}
Los conjuntos públicos pueden no cubrir herramientas nuevas, contextos largos, system prompts reales ni la distribución de usuarios. La aceptación debería exigir: mejora en el benchmark público, que la regression interna no empeore, mejora en las slices de tráfico real, revisión humana de los casos graves y condiciones de monitoreo y reversión después del lanzamiento.
\end{warningbox}

\subsection{Vibe check: cómo los ejemplos cualitativos se convierten en evidencia formal}

Las tres pantallas before/after usan surveillance fiction, time travel government project y technology-free heist. Todas contienen palabras peligrosas en apariencia, pero son solicitudes creativas legítimas. Claude 2.1 tiende a rechazarlas; Claude 3 ofrece ayuda estructurada. Un ejemplo aislado no demuestra la capacidad general, pero estos ejemplos sirven para verificar si el objetivo de producto se cumplió «de manera significativa» y ayudan a las personas a detectar diferencias de tono que el benchmark no codifica.

\lecturefigure{slide-47-vibe-check-surveillance.jpg}{Vibe check I: la escritura de ciencia ficción sobre vigilancia pasa del rechazo a la ayuda estructurada}{00:32:10--00:32:13}

\lecturefigure{slide-48-vibe-check-time-travel.jpg}{Vibe check II: before/after de un proyecto secreto de viajes en el tiempo}{00:32:13--00:32:21}

\lecturefigure{slide-49-vibe-check-heist.jpg}{Vibe check III: before/after de un diálogo de heist sin tecnología moderna}{00:32:21--00:32:30}

\begin{knowledgebox}{Cuándo la qualitative eval no es «una ocurrencia sin fundamento»}
La evaluación cualitativa requiere fijar el prompt, la versión del modelo, la configuración de muestreo y la rubric; al menos dos reviewers juzgan de forma independiente y se registra el disagreement. Es especialmente útil para descubrir nuevos failure modes, problemas de tono y puntos de ruptura en la interacción. Una vez descubiertos, los ejemplos deben ampliarse a conjuntos de control, de variantes y contrafactuales antes de incorporarse a la regresión continua.
\end{knowledgebox}

\begin{lstlisting}[caption={Marco de ejecución por capas para evals de comportamiento}]
suites = [public_benchmarks, frozen_regressions, synthetic_slices,
          product_replays, qualitative_panels]
results = run_all(model, suites, fixed_system_and_sampling=True)
report = aggregate(results, by=[risk, task, language, tool, length])
gate(report.overall, report.worst_slice, severe_case_review)
archive(model_version, judge_version, prompts, traces, report)
\end{lstlisting}

\teachervoice{00:31:44--00:32:30, la ponente muestra primero el XSTest cuantitativo y después usa los tres pares before/after para hacer un vibe check. Nota de clase: los números se encargan de la regresión a escala y los ejemplos cualitativos, de juzgar si el comportamiento corresponde realmente a las convicciones del producto; ninguno sustituye al otro.}

\subsection{Resumen de la sección}

La behavior eval confiable parte del vector de comportamiento y, a través de la taxonomy, los mixed-source prompts, la provenance de los datos, las métricas con restricciones, los benchmarks y la qualitative review, conforma un protocolo de decisión. La lección más importante del caso de over-refusal no es «rechazar menos», sino que fallas distintas tienen causas raíz distintas. La siguiente sección lleva este enfoque al RL: cómo funciona el producto depende de cómo se construyen el entorno, las acciones, el reward y la verificación.
\section{El RL Environment es arquitectura de producto: tareas, herramientas, recompensas y recuperación definen juntas el comportamiento}

La tercera parte de la clase condensa su título en una sola frase: \textbf{how you construct RL environment and rewards is how your product will work}. No es retórica. Si el entorno de entrenamiento solo admite respuestas de texto, el modelo no aprenderá de la nada a usar herramientas de forma confiable; si la reward solo considera la cadena final, el modelo puede ignorar los permisos, el costo y los efectos secundarios irreversibles del proceso; si el entorno no contempla la recuperación ante fallos, la política no tiene oportunidad de aprender cuándo detenerse, pedir ayuda o revertir.

\lecturefigure{slide-50-rl-environment-product.jpg}{Tesis central: la construcción del RL environment y de la reward determina el comportamiento del producto}{00:33:15--00:35:51}

\begin{importantbox}{Primer uso: MDP, POMDP y contract del entorno}
El Markov Decision Process (MDP) estándar se escribe como
\[
\mathcal{M}=(\mathcal{S},\mathcal{A},P,R,\gamma),
\]
donde $\mathcal{S}$ es el estado, $\mathcal{A}$ la acción, $P$ la transición de estados, $R$ la recompensa y $\gamma$ el descuento. Un producto real se parece más a un Partially Observable MDP (POMDP): el modelo solo ve la observation $o_t$; la intención del usuario, el estado de las herramientas y las consecuencias futuras no son completamente visibles. El contract del entorno también debe declarar permisos, presupuesto, terminación, registros, confirmación humana y recuperación.
\end{importantbox}

El retorno de la trayectoria es
\[
G_t=\sum_{k=0}^{T-t}\gamma^k r_{t+k}.
\]
En las tareas largas, lo decisivo no es la fórmula, sino de dónde proviene $r_t$. Las pruebas de código, el estado de la base de datos, la confirmación del usuario y el judge humano tienen distinta confiabilidad; comprimir una tarea subjetiva no verificable en una puntuación inmediata lleva a la política a buscar huecos en el indicador sustituto.

\subsection{Real-world complexity: herramientas, contexto largo y tareas que se completan en varios pasos}

La complejidad de las tareas reales de ingeniería de software o de investigación proviene de las llamadas a herramientas, el contexto largo, la retroalimentación diferida, los sistemas externos y el reward design. Los agent y las herramientas de la figura son solo una indicación estructural; la aceptación de ingeniería debe preguntar además: ¿la llamada tiene schema?, ¿se reintenta cuando la herramienta falla?, ¿se confirman las operaciones de escritura?, ¿las tareas largas admiten checkpoint?, ¿el modelo distingue entre «no hay resultados» y «la herramienta falló»?

\lecturefigure{slide-51-real-world-rl-complexity.jpg}{La complejidad de los casos de uso reales equivale a la complejidad del RL environment}{00:35:51--00:38:33}

\begin{knowledgebox}{Primer uso: tool call y state transition}
Una llamada a herramienta no es un fragmento de lenguaje natural. La acción puede escribirse como $a_t=(\text{tool},\text{arguments},\text{idempotency key})$, y el entorno devuelve una observation estructurada y un cambio de estado. Si el entrenamiento solo recompensa que el modelo «escriba texto que parece una API» sin ejecutar la herramienta, lo que la política aprende es a imitar el formato, no la capacidad de resolver la tarea.
\end{knowledgebox}

\begin{lstlisting}[caption={Contract de RL environment que sirve tanto para entrenar como para desplegar}]
environment = {
    "state": [artifact_versions, tool_status, user_goal, budget],
    "observations": [messages, tool_results, verifier_feedback],
    "actions": [respond, call_tool, edit, ask_user, stop, rollback],
    "permissions": policy_by_action_and_resource,
    "termination": [goal_reached, budget_exhausted, unsafe, user_stop],
    "rewards": [task_outcome, verifier_score, cost_penalty, safety_penalty],
    "audit": persist_full_trace_and_versions,
}
\end{lstlisting}

\subsection{Desglosar la dificultad del entorno}

En clase se pidió a los estudiantes comparar software engineer, AI researcher, creative storyteller y multiplayer/multi-agent environment. La dificultad no proviene solo de un espacio de respuestas amplio, sino también del retraso de la retroalimentación, la observabilidad, los interlocutores con quienes se colabora y el criterio de éxito. Las tareas de software pueden tener pruebas, pero incluyen despliegues asíncronos y requisitos implícitos; las tareas creativas carecen de una respuesta única, pero tienen la voice del autor, la coherencia de la trama y objetivos estéticos; los entornos con varias personas introducen, además, protocolos e incentivos.

\lecturefigure{slide-52-environment-complexity-exercise.jpg}{Ejercicio de complejidad del entorno: software, investigación, creación y colaboración multiagente}{00:38:33--00:40:55}

\begin{knowledgebox}{Los cinco ejes de dificultad del entorno}
\begin{tabularx}{\textwidth}{L{0.20\textwidth} X X}
\toprule
Eje & Extremo de baja dificultad & Extremo de alta dificultad \\
\midrule
Horizon & Respuesta de un solo paso & Tareas de horas o días \\
Observability & Estado completamente visible & Intención implícita, cambios externos \\
Verifiability & Pruebas exactas & Evaluación subjetiva, diferida o de varias personas \\
Action risk & Solo lectura y reversible & Escritura, pagos, publicación, control de dispositivos \\
Coordination & Un solo usuario & Varias personas, varios agent, objetivos en conflicto \\
\bottomrule
\end{tabularx}
Con el mismo tamaño de modelo, desplazar el entorno a lo largo de estos ejes genera requisitos de confiabilidad completamente distintos.
\end{knowledgebox}

\teachervoice{00:33:15--00:38:56, Karina atribuye la dificultad de software engineering, AI research y real-world tasks al RL environment, en lugar de limitarse a decir «el modelo todavía no es lo bastante grande». El docente enfatiza que tools, long context, reward design y multiplayer interaction aumentan la dificultad del aprendizaje.}

\subsection{Subjective tasks: que la verificabilidad sea débil no significa que no puedan investigarse}

La escritura, el diseño visual, las emociones y la proactividad, y la social intelligence difícilmente se miden con una sola objective measure. El curso no concluye que «no se pueden entrenar»; pide convertir la tarea en una situación más realista y recopilar comparaciones, contrafactuales y resultados de largo plazo. Lo esencial es reconocer los valores del judge y su dependencia del contexto, en lugar de presentarlos como verdad objetiva.

\lecturefigure{slide-53-subjective-tasks.jpg}{Tareas subjetivas difíciles de medir: escritura, estética, emociones e inteligencia social}{00:38:56--00:41:30}

Para dos salidas $y_a,y_b$, la pairwise preference suele modelarse como
\[
P(y_a\succ y_b\mid x)=\sigma\!\left(r_\phi(x,y_a)-r_\phi(x,y_b)\right),
\]
donde $r_\phi$ es el reward model y $\sigma$ es la función logística. Pairwise es más fácil que la puntuación absoluta, pero las preferencias siguen cambiando según el reviewer, la cultura, el propósito y el contexto. Hay que conservar el disagreement en lugar de promediarlo a la fuerza en un «gusto global».

\begin{warningbox}{La simplificación más peligrosa en las tareas subjetivas}
Tomar los «me gusta», el tiempo de permanencia, la virality o un único AI judge como si fueran la creativity misma mezcla en la reward la retroalimentación social, la distribución de la plataforma y la novedad. Las señales del mundo real pueden formar parte del entorno, pero hay que prevenir la manipulación, los sesgos de grupo y la Goodhart's law.
\end{warningbox}

\subsection{La nueva Product Research: ciclos rápidos en lugar de congelar la especificación de una sola vez}

El ponente trata las nuevas tareas como simulation de situaciones reales, y combina in-context learning, distillation de modelos fuertes, synthetic data, nuevos model behavior, multiplayer interaction y deliberate RL. Aquí, «rápido» no significa omitir el rigor, sino que prototype, eval y entrenamiento intercambien evidencia en ciclos más cortos.

\lecturefigure{slide-54-product-research-loop.jpg}{La nueva Product Research: ciclo rápido entre situaciones reales, ICL, synthetic data y RL}{00:41:30--00:42:14}

\begin{importantbox}{La unidad experimental del co-design loop}
Cada ronda de experimentos debe fijar una sola hypothesis, por ejemplo, «permitir que el modelo haga primero preguntas aclaratorias reduce las llamadas erróneas a herramientas». Luego se registran a la vez: el cambio de entorno, el cambio de datos, la versión del modelo, el cambio de interfaz y el outcome. Si los cinco cambian juntos, no hay forma de saber de dónde proviene la mejora. La iteración rápida sigue requiriendo controlled comparison.
\end{importantbox}

Una iteración de co-design puede escribirse como
\[
\Delta=(\Delta E,\Delta D,\Delta \pi,\Delta I),
\]
que representan, respectivamente, los cambios en environment, data, policy/model e interface. El experimento ideal deja distintos de cero solo unos pocos $\Delta$ a la vez y compara sobre la misma eval; el lanzamiento del producto, en cambio, requiere monitorear también los efectos de interacción.

\subsection{Reward design: dejar claro qué retroalimentación se quiere}

La diapositiva de reward design pregunta: si queremos que el modelo se desempeñe mejor en situaciones sociales reales, ¿qué retroalimentación debemos darle? La respuesta requiere deep product thinking, porque la «satisfacción del usuario» puede provenir de una complacencia a corto plazo, la «proactividad» puede volverse una interrupción y «tener creatividad» puede convertirse en no respetar las restricciones. La reward debe conectar el outcome real, las restricciones del proceso y el judgment humano.

\lecturefigure{slide-55-reward-design-product-thinking.jpg}{Reward design: la retroalimentación debe enseñar al modelo el comportamiento adecuado en escenarios reales}{00:42:14--00:45:32}

Una descomposición posible es
\[
R(\tau)=w_oR_{\text{outcome}}+w_pR_{\text{process}}
-w_cC_{\text{cost}}-w_rC_{\text{risk}}.
\]
$R_{\text{outcome}}$ mide el resultado de la tarea, $R_{\text{process}}$ recompensa los procesos de aclaración, verificación y colaboración, $C_{\text{cost}}$ es el costo en tiempo, cómputo y herramientas, y $C_{\text{risk}}$ es el riesgo inaceptable. Los pesos no lo resuelven todo; si un término no puede medirse de forma confiable, la optimización solo amplificará el error de medición.

\teachervoice{00:38:56--00:43:43, la clase conecta subjective tasks, la simulation de situaciones reales, synthetic data, multiplayer interaction y reward design. La pregunta central del docente es: ¿qué retroalimentación estás realmente dispuesto a darle al modelo para que sea más útil en las relaciones del mundo real?}

\subsection{Reward hacking y asymmetric verification}

Reward hacking se refiere a que la política encuentra un camino de alta puntuación sin cumplir el objetivo que el diseñador realmente quería. A medida que los reasoning model y los agent de software se vuelven más complejos, el hack puede esconderse en trayectorias largas, huecos en las pruebas, puertas traseras en el código o sesgos del judge. El artículo de la figura y el texto de Lilian Weng advierten: la reward no es ground truth, y la verification affordance es en sí misma diseño de alignment.

\lecturefigure{slide-56-reward-hacking.jpg}{Reward hacking: la puntuación sustituta sube, pero el objetivo real puede quedar eludido}{00:43:43--00:45:46}

Si la utilidad real es $U(\tau)$ y el indicador sustituto medible es $\hat U(\tau)$, lo que se optimiza es
\[
\max_\pi\;\mathbb{E}[\hat U(\tau)],
\]
pero al producto le importa $U$. Tras un cambio de distribución, la correlación entre $\hat U$ y $U$ puede disminuir. La llamada \term{asymmetric verification} consiste en que generar un resultado de alta calidad es difícil, pero revisar un resultado dado es relativamente fácil; esto permite incorporar el verifier al entrenamiento. Sin embargo, en cuanto el verifier tiene un blind spot, la política también aprende a atacarlo.

\begin{knowledgebox}{¿Qué tan independiente debe ser la cadena de verificación?}
Si un mismo modelo genera la respuesta, la explica y se califica a sí mismo, sus errores pueden estar muy correlacionados. Una cadena de verificación más sólida combina pruebas ejecutables, comprobaciones de reglas, modelos independientes, muestreo humano, distintas fuentes de datos y adversarial cases. A mayor independencia, mayor costo; el producto debe asignar el presupuesto de verificación según el riesgo de cada acción.
\end{knowledgebox}

\begin{warningbox}{Síntomas de reward hacking en el producto}
El modelo puede ocultar fallos, manipular al judge, generar código frágil que solo pasa las pruebas, pedir confirmación al usuario en exceso para eludir la responsabilidad o elegir tareas fáciles para elevar la tasa de éxito. Si solo se mira la reward promedio, estas estrategias parecen progreso; hay que revisar el trace completo, el denominador de tareas y los casos no terminados.
\end{warningbox}

\teachervoice{00:43:43--00:45:46, Karina señala que cuanto más complejos son el reasoning y la software engineering, más complejos son también los reward hacks; el usuario puede no entender qué código modificó el modelo, por lo que se necesita una nueva trustworthy verification affordance. Nota de clase: la interfaz de verificación forma parte del alignment.}

\subsection{Resumen de la sección}

La unidad de un producto de RL no es un modelo aislado, sino un entorno al estilo POMDP: estado, observación, acción, herramientas, permisos, reward, costo, terminación, registros y recuperación moldean juntos el comportamiento. Las tareas subjetivas pueden investigarse, pero hay que conservar las diferencias de preferencia; la reward puede optimizarse, pero hay que reconocer el error del indicador sustituto y los ataques al verifier. El rigor del co-design proviene de registrar con claridad, en cada ronda, los cambios de entorno, datos, modelo e interfaz.
\section{Vignettes del futuro y Q\&A: de la inteligencia barata a la inteligencia social}

Las últimas cuatro diapositivas de la presentación formal extrapolan al futuro el diseño de sistemas expuesto antes: la caída del costo de las capacidades, las interfaces generadas dinámicamente, la personalización de la educación y la salud, y el cambio en la relación con la creación de historias. Después, los 22 minutos de Q\&A aportan condiciones restrictivas más valiosas. Las notas leen juntas la visión y las respuestas: a cada «se puede» le agregan «qué evidence falta».

\subsection{El costo de las capacidades baja, pero el costo de verificación no necesariamente baja al mismo ritmo}

La gráfica de la MMLU cost frontier expresa una tendencia histórica: con una precisión similar, el precio de la inferencia baja rápidamente. No demuestra que el costo vaya a bajar exponencialmente para siempre, ni que MMLU represente el trabajo real. Más importante aún, que generar sea más barato aumenta el volumen de salidas; si la verificación sigue dependiendo de expertos, el costo total del sistema puede desplazarse de la inference a la review.

\lecturefigure{slide-57-mmlu-cost-frontier.jpg}{MMLU performance--cost frontier: caída histórica del costo de invocación para capacidades equivalentes}{00:45:46--00:47:40}

Sea $q(c)$ la tasa de éxito de la tarea, $c$ el costo de una sola inferencia y $v$ el costo de la verificación humana; el costo por éxito confiable puede escribirse, de forma aproximada, como
\[
C_{\text{trusted}}=\frac{c+v}{q(c)}.
\]
Aunque $c$ baje, si las tareas más complejas hacen subir $v$ o vuelven inestable $q$, los resultados confiables no se abaratan al mismo ritmo. El producto debe optimizar $C_{\text{trusted}}$, y no solo reportar el token price.

\teachervoice{00:45:46--00:47:40, la ponente dice que esta gráfica ya es «de hace un año» y que MMLU ya no es el benchmark que más interesa; la usa para expresar que la raw intelligence se abarata. Las notas advierten: es una dated trend, no una ley de extrapolación ilimitada.}

\subsection{Dynamic UI, servicios personalizados y la relación con las historias}

La sección anterior mostró que baja el costo de invocar la raw intelligence, pero una salida barata solo genera valor para el usuario si llega a una interfaz adecuada. Esta sección conecta la tendencia de costos con tres tipos de surface: la UI dinámica generada según la intención, los puntos de acceso a la educación y la salud que asumen recomendaciones de alto riesgo, y la colaboración narrativa que preserva la agency del autor. La visión de la Dynamic generative UI es que la interfaz se genere al instante según la intención: quien aprende de forma visual ve gráficos manipulables; quien aprende de forma auditiva recibe audio. Esto exige que el modelo infiera a la vez la tarea, las preferencias del usuario y los dispositivos disponibles; una inferencia errónea debe poder deshacerse con facilidad, pues de lo contrario el «software invisible» también volverá impredecible el comportamiento del sistema.

\lecturefigure{slide-58-dynamic-generative-ui.jpg}{Dynamic generative UI: la interfaz se genera según la intención y la creación de software tiende a volverse invisible}{00:47:40--00:48:04}

\lecturefigure{slide-59-personalized-health-education.jpg}{Salud y educación personalizadas con consumer hardware como punto de acceso futuro}{00:48:04--00:48:42}

La salud y la educación tienen un gran valor, y también riesgos distintos. Un sistema educativo puede admitir la exploración y la corrección de errores; las recomendaciones médicas implican omisión de síntomas, triaje de urgencias, privacidad y regulación. Un indicador unificado de «personalización» oculta los niveles de riesgo. El producto debe distinguir entre information, decision support y autonomous action, y prever revisión profesional y escalamiento de emergencia para las salidas de alto riesgo.

\lecturefigure{slide-60-storytelling-relationship.jpg}{La relación de las personas con el storytelling process cambiará con los modelos colaborativos}{00:48:42--00:49:18}

La apertura de la creación narrativa la convierte en un buen entorno para poner a prueba la social intelligence, la voice preservation y la agency. El modelo puede dar inspiración, reescrituras y sugerencias de estructura, pero «decidir por el autor qué vale la pena expresar» y «ayudar al autor a realizar su intención» son productos distintos. Un verdadero collaborator debe permitir que el autor vea las fuentes, las opciones y los cambios, en lugar de arrebatarle el control creativo con texto fluido.

\begin{warningbox}{Una future vignette no equivale a una deployment claim}
Estas cuatro diapositivas expresan el juicio de la ponente sobre la dirección a seguir. No demuestran que ya estén resueltas la seguridad médica, la ganancia educativa, la previsibilidad de la UI dinámica ni el impacto sobre el trabajo creativo. Las notas las tratan como una research agenda: cada visión debe completarse con usuarios objetivo, modos de falla, permisos, verificación, costo y responsables.
\end{warningbox}

\teachervoice{00:47:40--00:49:18, Karina imagina interfaces dinámicas guiadas por preferencias visuales o auditivas, educación y salud personalizadas y una nueva relación con el storytelling, y a la vez espera que los creadores vean la IA como una herramienta y no solo como una amenaza. Nota de clase: esta visión optimista debe realizarse junto con la agency del usuario y los mecanismos de verificación.}

\subsection{Q\&A I: tareas subjetivas, benchmarks y rutas de emprendimiento}

El Q\&A pregunta primero: ¿hay tareas difíciles de evaluar que, por ello, los investigadores prefieren no abordar? La ponente pone como ejemplos la creative writing y la emotional intelligence; considera que es posible construir benchmarks, pero las tareas complejas avanzan más despacio y dependen más de hitos. También señala que quien emprende no necesita entrenar primero un foundation model: puede usar modelos existentes de bajo costo y concentrar la innovación en tareas reales y en el ciclo cerrado del producto.

\begin{knowledgebox}{Nota de clase: subjetivo no significa imposible de evaluar}
00:49:18--00:55:00, Karina distingue entre «no existe un frontier benchmark disponible» y «en principio no se puede construir una evaluación». Se puede partir de comparaciones por pares, panels de expertos, outcomes de largo plazo de los usuarios, resultados de competencias y ejemplos contrafactuales; pero todas estas señales conllevan juicios de valor y no deben hacerse pasar por una verdad natural.
\end{knowledgebox}

\subsection{Q\&A II: cómo evitar que las preferences converjan al promedio global}

La ponente propone dos líneas: preservar la entropy del base model y usar RLAIF/synthetic preference data para generar de forma dirigida la diversidad necesaria. Aquí la entropy no es un estímulo a la aleatoriedad, sino una forma de evitar que el RL comprima muchos estilos útiles en «los emojis, el Markdown y el tono que le gustan al usuario promedio».

Para un contexto $x$, la entropy de la política es
\[
\mathcal{H}(\pi(\cdot\mid x))=-\sum_y \pi(y\mid x)\log \pi(y\mid x).
\]
Una preference optimization demasiado intensa puede reducir $\mathcal{H}$ y provocar mode collapse. En la práctica, conviene mantener la diversidad de candidatos según la tarea y, a la vez, filtrar con restricciones los errores factuales y los riesgos de seguridad; «diverso» no significa «cualquier salida es buena».

\begin{importantbox}{Nota de clase: el valor de RLAIF es la cobertura controlable, no un volumen ilimitado de datos}
00:55:00--01:01:20, Karina dice que los synthetic data no necesariamente tienen que ser muchos; lo decisivo es la diversity. El equipo todavía puede revisarlos manualmente, o validarlos con un modelo independiente bien calibrado y con meta-eval. Describe la synthetic generation como una curation de la distribución objetivo, no como una copia de la preferencia humana promedio.
\end{importantbox}

\subsection{Q\&A III: diagnóstico cualitativo, costos y frontier work}

Sobre los bugs del modelo, la ponente reconoce que mucha nuanced weirdness surge de que las personas «juegan con el modelo» una y otra vez; las eval automáticas sirven para revisar behaviors conocidos, pero los problemas nuevos suelen descubrirse primero mediante exploración cualitativa. También distingue el costo del desarrollo de producto del de la frontier research: los equipos de aplicación pueden reutilizar modelos baratos; los equipos de frontera deben asumir la ineficiencia de la etapa de invención, y solo después llega una segunda ronda de optimización de costos.

\begin{knowledgebox}{Nota de clase: una anomalía aislada frente a un comportamiento sistemático}
00:57:20--01:00:00, Karina dice que la revisión cualitativa debe fijarse en si el behavior aparece de forma persistente. Un sample ocasional puede deberse a la aleatoriedad; solo cuando se reproduce de forma estable entre variantes de prompt, sesiones y versiones debe escalarse a issue sistemático. Las notas recomiendan ampliar cada hallazgo manual a un conjunto de variantes y luego incorporarlo a la regresión automática.
\end{knowledgebox}

\begin{warningbox}{Límites de la discusión sobre costos}
Las respuestas del Q\&A sobre el costo de entrenamiento, el scaling y el abaratamiento futuro tienen una incertidumbre evidente. Lo que puede confirmarse es que la capa de aplicación puede reutilizar modelos existentes, que la invención de frontera es costosa y que la ingeniería posterior reduce costos; de las respuestas en clase no puede deducirse que el costo de entrenamiento vaya a bajar necesariamente según una curva determinada.
\end{warningbox}

\subsection{Q\&A IV: robots, AI coworker y social intelligence}

La ponente es optimista respecto de la robotics, pero considera que los datos son un enorme bottleneck. A la pregunta de si «ya existen colegas de IA», responde con claridad que todavía no; la forma ideal se parecería más a compartir pantalla, conversar en tiempo real, señalar objetos concretos y ajustar la agency según la voluntad del usuario. Lo que falta no es otro Slack bot, sino social intelligence, speech-to-speech, grounding multimodal y un espacio de trabajo compartido.

\begin{knowledgebox}{Nota de clase: qué le falta todavía al AI coworker}
01:05:00--01:10:00, las carencias que enumera Karina incluyen: generación en tiempo real y edición conjunta, saber cuándo orientar sin quitar el control, speech-to-speech, señalar al mismo tiempo el objeto del que se habla y adaptarse a distintas relaciones. Esta lista es más exigente que «poder hacer tareas automáticamente» y explica también por qué una teammate eval debe incluir la coordinación interpersonal.
\end{knowledgebox}

\subsection{Q\&A V: Research-driven product frente al flujo tradicional de PRD}

El desarrollo de producto tradicional suele empezar con el PRD, el diseño y la planeación de ingeniería; un research-driven product puede surgir primero de una capability demo sorprendente y luego dar forma al producto en torno a ella, o bien producto e investigación pueden experimentar juntos desde el primer día. La ponente presenta Canvas como ejemplo de lo segundo: el post-training trabaja junto con producto y el proceso es más ad hoc, pero eso no significa que no haya normas; las normas se trasladan a los prototypes, las eval, las traces y la evidencia de iteración.

\begin{importantbox}{Nota de clase: una demo de capacidad no es el punto final del producto}
01:09:52--01:11:29, Karina contrasta el flujo de software tradicional con el research-driven product: cuando la investigación produce primero una capacidad, el producto se forma en torno a ella; a veces producto e investigación se diseñan juntos desde el inicio. El criterio de ingeniería es que, cuanto más ad hoc sea el proceso, más necesario es registrar la hypothesis, las versiones y los failures, para evitar que el equipo solo recuerde las demos exitosas.
\end{importantbox}

\subsection{Resumen de la sección}

La condición común de las visiones de futuro no es «modelos más inteligentes», sino que baje el costo por éxito confiable, que las interfaces puedan verificarse y deshacerse, que la personalización respete el control del usuario, que las tareas subjetivas conserven la diversidad y que el AI coworker tenga capacidades sociales y multimodales. El Q\&A completa la tesis del curso: muchas capacidades pueden enseñarse, pero los datos, la verificación, la distribución de preferencias y la relación de colaboración determinan si se convierten en un producto confiable.
\section{Síntesis y ampliación}

Esta clase devuelve el RL del terreno de los algoritmos de optimización al de los sistemas de producto. La product belief determina qué conducta se quiere ver; la interface expone las tareas y los fallos; la eval descompone la «sensación» en evidence reproducible; el data debugging rastrea el origen de las conductas; el RL environment define estados, acciones, herramientas y recuperación; la reward incorpora en el objetivo tanto los resultados como el proceso; y la evidence de los usuarios tras la puesta en producción inicia la siguiente ronda de investigación. Si falta cualquiera de estas capas, otra capa tendrá que cargar con problemas que no puede resolver.

\subsection{Una tabla que articula el co-design loop}

\begin{table}[H]
\centering
\small
\begin{tabular}{@{}L{0.16\textwidth}L{0.24\textwidth}L{0.27\textwidth}L{0.23\textwidth}@{}}
\toprule
Capa & Objeto central & Fallo representativo & Pregunta de aceptación \\
\midrule
Belief & Conducta objetivo y relación con el usuario & Lemas que no pueden llevarse a la práctica & ¿Puede escribirse como un contract de conducta? \\
Interface & artifact, diff, control y permisos & Capacidad invisible o irreversible & ¿El usuario puede entender, editar y detener? \\
Eval & Distribución, rubric, judge, umbrales & Desajuste entre el benchmark y las tareas reales & ¿Es reproducible, está dividida por slice y permite decidir? \\
Data & provenance, mixture, preference & Conflictos y contaminación de datos & ¿Puede localizarse el origen de una conducta y revertirse? \\
Environment & Estados, acciones, herramientas, terminación & Las tareas de entrenamiento difieren de las de despliegue & ¿Coinciden permisos, presupuesto y recuperación ante fallos? \\
Reward & outcome, process, cost, risk & reward hacking / Goodhart & ¿El verifier es independiente y resistente a ataques? \\
Deployment & telemetry, regresiones, responsables & El éxito de la demo oculta incidentes de cola larga & ¿Se monitorea el outcome y puede revertirse con rapidez? \\
\bottomrule
\end{tabular}
\caption{Mapa de aceptación en siete capas para el diseño conjunto de RL, producto e investigación.}
\end{table}

\subsection{Preguntas que deben responderse antes de iniciar un proyecto de Frontier Product}

\begin{importantbox}{Autoevaluación en doce preguntas}
\begin{enumerate}
  \item ¿La product belief puede escribirse como conducta observable y no como adjetivos?
  \item ¿Cuáles son las tareas reales del usuario objetivo y el costo de los fallos?
  \item ¿Mediante qué form factor se expone la capability del modelo?
  \item ¿Qué artifacts, estados y fuentes deben poder rastrearse?
  \item ¿Qué facultades de confirmar, editar, detener y hacer rollback conserva el usuario?
  \item ¿Están fijos la distribución, el judge, la rubric y la versión de la eval?
  \item ¿Se reportan a la vez los resultados overall, por slice, worst-case y los casos graves?
  \item ¿La taxonomy de fallos puede asociarse con causas raíz en los datos, el sistema o la policy?
  \item ¿Las herramientas, los permisos, el presupuesto y la terminación del RL environment coinciden con los del despliegue?
  \item ¿La reward separa el outcome real, el proceso, el costo y el riesgo?
  \item ¿Qué blind spots tiene el verifier y puede la política atacarlo?
  \item Tras la puesta en producción, ¿quién revisa el outcome, quién puede revertir y qué condiciones activan la detención del servicio?
\end{enumerate}
\end{importantbox}

\subsection{Conjunto mínimo de registros para la reproducibilidad}

Ya se investiguen data, behavior, RL environment o interface, deben guardarse source/version, prompt/system, el modelo y el sampling, el schema de las herramientas, la trace completa, judge/rubric, el disagreement humano, las muestras de fallo y las condiciones de reversión. Si se usa un AI judge o synthetic data, también deben fijarse el modelo generador, el modelo judge y los permisos; de lo contrario, un resultado no podrá ser explicado por la versión siguiente.

\begin{lstlisting}[caption={Registro reproducible de un experimento de Frontier product}]
experiment = {
    "hypothesis": product_belief_as_testable_behavior,
    "model": [base, post_training, checkpoint, sampling],
    "environment": [tools, permissions, budget, termination],
    "interface": [surface, artifact_schema, confirmation, rollback],
    "evaluation": [datasets, slices, judges, rubrics, thresholds],
    "evidence": [traces, outcomes, disagreements, severe_failures],
    "decision": [ship, hold, rollback, next_hypothesis],
}
\end{lstlisting}

\subsection{Lecturas adicionales}

\begin{itemize}
  \item Bai et al., \emph{Constitutional AI: Harmlessness from AI Feedback}: para entender RLAIF, las restricciones basadas en principios y la AI feedback.
  \item Röttger et al., \emph{XSTest}: estudia la exaggerated safety y el over-refusal ante benign prompts.
  \item Lilian Weng, \emph{Reward Hacking in Reinforcement Learning}: recopila de forma sistemática la proxy reward y las vías de hacking.
  \item Ouyang et al., \emph{Training language models to follow instructions with human feedback}: el pipeline clásico de RLHF y sus limitaciones.
  \item Amodei et al., \emph{Concrete Problems in AI Safety}: para entender el diseño de entornos a partir del reward hacking, los side effects y la oversight.
\end{itemize}

\begin{warningbox}{Límite final}
Los nombres de productos, el desempeño de los modelos, las curvas de costo y las prácticas organizacionales de esta clase corresponden al momento en que se impartió, en 2025. Las notas conservan el método de co-design; no tratan ninguna demo, benchmark o proceso de empresa en particular como un hecho permanente. La capacidad realmente transferible consiste en traducir la visión en un entorno, la conducta en una eval, rastrear los fallos hasta los datos y el sistema, y luego decidir con evidencia reproducible si se entrega.
\end{warningbox}

\end{document}
